\section{Параметры}

\begin{table}[htbp]
	\centering
	\small
	\begin{tabularx}{\textwidth}{l c c X}
		\toprule
		\textbf{Параметр} & \textbf{Допустимые значения} & \textbf{По умолчанию} & \textbf{Описание} \\
		\midrule
		
		\verb|ADDR_WIDTH| &
		Положительное целое &
		\verb|32| &
		Разрядность адреса памяти. \\
		
		\verb|DATA_WIDTH| &
		32, 64 &
		\verb|32| &
		Разрядность шины данных. \\
		
		\verb|SET_NUM| &
		Степень двойки &
		\verb|4| &
		Число наборов кэша. \\
		
		\verb|WAY_NUM| &
		Степень двойки &
		\verb|4| &
		Число секций кэша (степень ассоциативности). \\

		\verb|LFSR_WIDTH| &
		Положительное целое &
		--- &
		Число триггеров в сдвиговом регистре \PRG \\
		
		\verb|LFSR_POLY| &
		Настраивается &
		--- &
		Полином обратной связи LFSR в \PRG. \\
		
		\verb|DEBUG_ON| &
		0, 1 &
		\verb|0| &
		Включает \StatisticsBlock при значении \verb|1|. \\
		
		\bottomrule
	\end{tabularx}
	\caption{Параметры настройки RTL}
	\label{tab:rtl_parameters}
\end{table}

Вспомогательные параметры ниже обозначают разрядности индексов в документации и не являются параметрами RTL.

\begin{table}[htbp]
	\centering
	\small
	\begin{tabularx}{\textwidth}{l c X}
		\toprule
		\textbf{Параметр} & \textbf{Определение} & \textbf{Описание} \\
		\midrule
		
		\verb|SET_WIDTH| &
		$\log_2(\text{\verb|SET_NUM|})$ &
		Разрядность индекса набора кэша. \\
		
		\verb|WAY_WIDTH| &
		$\log_2(\text{\verb|WAY_NUM|})$ &
		Разрядность индекса секции. \\
		
		\bottomrule
	\end{tabularx}
	\caption{Вспомогательные параметры документации}
	\label{tab:documentation_parameters}
\end{table}
